%% Research manuscript — with inline figures
%% Compile: xelatex manuscript && xelatex manuscript

\documentclass[10pt,twocolumn]{article}

% ── Packages ───────────────────────────────────────────────────────────────
\usepackage{fontspec}
\usepackage[margin=1.8cm, top=2cm, bottom=2.2cm,
            columnsep=0.6cm]{geometry}
\usepackage{microtype}
\usepackage{graphicx}
\usepackage{booktabs}
\usepackage{tabularx}
\usepackage{array}
\usepackage{xcolor}
\usepackage{amsmath}
\usepackage{amssymb}
\usepackage{url}
\usepackage{hyperref}
\usepackage{caption}
\usepackage{subcaption}
\usepackage{float}
\usepackage{titlesec}
\usepackage{parskip}
\usepackage{enumitem}
\usepackage{fancyhdr}
\usepackage{setspace}
\usepackage{authblk}
\usepackage{multirow}
\usepackage{siunitx}
\usepackage{etoolbox}
\usepackage{stfloats}   % enables figure* bottom placement

% ── Fonts ──────────────────────────────────────────────────────────────────
\setmainfont{Helvetica Neue}[
  BoldFont   = Helvetica Neue Bold,
  ItalicFont = Helvetica Neue Italic,
]
\setsansfont{Helvetica Neue}

% ── Colours ────────────────────────────────────────────────────────────────
\definecolor{natureblue}{RGB}{0, 80, 150}
\definecolor{naturered}{RGB}{200, 40, 40}
\definecolor{darkgray}{RGB}{80, 80, 80}
\definecolor{medgray}{RGB}{140, 140, 140}
\definecolor{accentgreen}{RGB}{30, 140, 90}
\definecolor{panelbg}{RGB}{248, 248, 248}

% ── Graphics path ──────────────────────────────────────────────────────────
\graphicspath{{../figures/}}

% ── Hyperref ───────────────────────────────────────────────────────────────
\hypersetup{
  colorlinks=true,
  linkcolor=natureblue,
  citecolor=natureblue,
  urlcolor=natureblue,
}

% ── Section formatting ─────────────────────────────────────────────────────
\titleformat{\section}
  {\normalfont\fontsize{9.5}{12}\bfseries\color{black}}{}{0em}{}
\titleformat{\subsection}
  {\normalfont\fontsize{9}{11}\bfseries\itshape\color{darkgray}}{}{0em}{}
\titlespacing*{\section}  {0pt}{12pt plus 2pt minus 2pt}{4pt}
\titlespacing*{\subsection}{0pt}{8pt  plus 2pt minus 1pt}{2pt}

% ── Caption styling ────────────────────────────────────────────────────────
\captionsetup{
  font={small},
  labelfont={bf},
  labelsep=period,
  format=plain,
  justification=justified,
  singlelinecheck=false,
}
\captionsetup[subfigure]{
  font={footnotesize},
  labelfont={bf},
  labelsep=space,
  labelformat=simple,
  justification=centering,
}

% ── Abstract env ───────────────────────────────────────────────────────────
\newenvironment{myabstract}{%
  \begin{minipage}{\linewidth}
  \rule{\linewidth}{1pt}\vspace{3pt}
  {\fontsize{8.5}{11}\selectfont\bfseries ABSTRACT\par}\vspace{3pt}
  \fontsize{8.5}{12}\selectfont
}{%
  \par\vspace{3pt}\rule{\linewidth}{0.5pt}
  \end{minipage}
}

% ── Header/footer ──────────────────────────────────────────────────────────
\pagestyle{fancy}
\fancyhf{}
\fancyhead[L]{\fontsize{7}{9}\selectfont\color{medgray}\textsc{Article}}
\fancyhead[R]{\fontsize{7}{9}\selectfont\color{medgray}\today}
\fancyfoot[C]{\fontsize{7}{9}\selectfont\color{medgray}\thepage}
\renewcommand{\headrulewidth}{0.3pt}
\renewcommand{\footrulewidth}{0pt}

% ── Misc ───────────────────────────────────────────────────────────────────
\setlength{\parindent}{0pt}
\setlength{\parskip}{5pt}
\setlist[itemize]{leftmargin=*, noitemsep, topsep=2pt}
\sisetup{separate-uncertainty, table-format=1.3}

\renewcommand\thesubfigure{\textbf{\alph{subfigure}}}

% ═══════════════════════════════════════════════════════════════════════════
\begin{document}

% ── Title block ────────────────────────────────────────────────────────────
\twocolumn[%
\begin{@twocolumnfalse}

  \vspace*{4pt}
  {\fontsize{7}{9}\selectfont\color{medgray}\textsc{Article}\hfill
   \textit{Received XX Month 2026\enspace|\enspace
           Accepted XX Month 2026\enspace|\enspace
           Published online XX Month 2026}\par}
  \vspace{8pt}

  {\fontsize{18}{22}\bfseries\raggedright
   p16/CDKN2A-mediated senescence shapes immune\par
   phenotype and patient survival in human cancer\par}

  \vspace{10pt}

  {\fontsize{8.5}{12}\selectfont
   David~C.~Kuo\textsuperscript{1,\dag}\par}

  \vspace{6pt}

  {\fontsize{7.5}{10}\selectfont\color{darkgray}
   \textsuperscript{1}Department of Computational Biology,
   [Institution], [City], [Country].\\
   \textsuperscript{\dag}Correspondence:
   \href{mailto:author@institution.edu}{author@institution.edu}\par}

  \vspace{10pt}

  \begin{myabstract}
Tumour immune phenotype — the distinction between immune-inflamed
(``hot'') and immune-excluded (``cold'') tumours — profoundly influences
immunotherapy response and patient outcomes. Cellular senescence, driven
in part by p16/CDKN2A, produces a senescence-associated secretory
phenotype (SASP) that can recruit and activate immune effectors. Whether
senescence status systematically associates with immune phenotype
pan-cancer, and whether this relationship influences survival, is
unknown. Here we performed an integrated analysis of RNA-seq, genomic
alteration, and clinical data from 9,910 primary tumour samples across 32 TCGA
cancer types. We show that CDKN2A expression positively correlates with
IFN-$\gamma$ score pan-cancer (pooled Spearman $\rho=0.158$; median
within-type $\rho=0.112$),
that SASP scores are enriched in hot tumours (Kruskal-Wallis $p<0.001$),
and that CDKN2A genomic deletion is enriched in cold tumours only after
stratifying by tumour type (Mantel-Haenszel OR\,=\,1.30, $p=4\times10^{-4}$;
the pooled $\chi^2$ test points the opposite way). A six-group
stratification (immune tertile\,$\times$\,p16-high/low) did not show the
pre-specified ordering, and its unstratified log-rank $p=3.3\times10^{-3}$
became non-significant when stratified by tumour type ($p=0.08$). Cox proportional hazards modelling
reveals IFN-$\gamma$ score as an independent survival predictor
(HR\,=\,0.907, 95\,\%\,CI\,0.856--0.961, $p=8.5\times10^{-4}$), while a
significant CDKN2A\,$\times$\,IFN-$\gamma$ interaction
(HR\,=\,0.964, $p=1.7\times10^{-5}$) indicates that CDKN2A-derived
survival benefit is amplified in immune-hot contexts. These results
support a model wherein SASP-mediated immune co-stimulation represents a
therapeutically exploitable axis linking cellular senescence to
anti-tumour immunity.
  \end{myabstract}

  \vspace{14pt}
\end{@twocolumnfalse}%
]

% ═══════════════════════════════════════════════════════════════════════════
\section{Introduction}

Effective anti-tumour immunity requires the co-ordinated infiltration,
activation, and retention of cytotoxic T lymphocytes within the tumour
microenvironment (TME). This requirement underpins the clinical
distinction between immune-inflamed (``hot'') tumours, characterised by
CD8$^+$ T-cell infiltration and IFN-$\gamma$ signalling, and
immune-excluded or immune-desert (``cold'') tumours, in which adaptive
immunity is dysfunctional or absent\cite{Chen2018}. The hot/cold
phenotype is a major determinant of immunotherapy response: checkpoint
blockade therapies produce durable responses in hot tumours at rates far
exceeding those seen in cold tumours, driving intense efforts to identify
factors that govern immune exclusion\cite{Sharma2017}.

Cellular senescence — a state of stable proliferative arrest typically
triggered by oncogene activation, genotoxic stress, or replicative
exhaustion — has emerged as a potential modulator of tumour
immunogenicity\cite{Coppé2010}. Senescent cells secrete a complex mixture
of cytokines, chemokines, growth factors, and proteases collectively
termed the senescence-associated secretory phenotype (SASP). SASP
components including IL-6, IL-8, CXCL1, CXCL10, and SERPINE1/PAI-1
have direct immunostimulatory activity: they recruit natural killer cells
and dendritic cells, activate NF-$\kappa$B, and can amplify interferon
signalling, thereby potentially promoting the transition from cold to hot
tumour states\cite{Acosta2013,Kuilman2008}.

The cyclin-dependent kinase inhibitor p16 (encoded by \textit{CDKN2A})
is a central mediator of the G1/S checkpoint and one of the most
frequently deleted tumour-suppressor genes in human
cancer\cite{Rayess2012}. At the same time, CDKN2A upregulation is a
reliable marker of senescence entry, co-ordinating SASP induction
through the p38-MAPK and NF-$\kappa$B pathways. Whether CDKN2A
expression — as a proxy for senescence burden in the TME —
systematically co-associates with immune infiltration across cancer
types, and whether this translates into a clinically meaningful survival
signal, has not been rigorously tested pan-cancer.

Here we address this gap using TCGA pan-cancer RNA-seq data from 32
primary-tumour types, integrated with genomic alteration calls, CIBERSORT immune
deconvolution, and clinical survival endpoints. We pre-specified five
confirmatory hypotheses (H1--H5) in the analysis notebooks (they were not
formally registered) and report their outcomes alongside
exploratory analyses, applying stringent Bonferroni correction
($\alpha=0.01$) to primary tests and Benjamini-Hochberg FDR
($\alpha=0.05$) for secondary analyses.

% ═══════════════════════════════════════════════════════════════════════════
\section{Results}

\subsection{Pan-cancer immune classification}

We classified 9,910 primary tumour samples across 32 TCGA cancer types
(acute myeloid leukaemia is absent because only primary-tumour samples
are retained) as hot, intermediate, or cold using a six-gene
IFN-$\gamma$ signature (\textit{IFNG}, \textit{STAT1}, \textit{CCL5},
\textit{CXCL9}, \textit{CXCL10}, \textit{IDO1}), applying a pan-cancer
tertile split (thresholds computed across all samples, not within
tumour type). The classification was moderately concordant with CD8$^+$
T-cell fraction estimates from Thorsson \textit{et al.}\cite{Thorsson2018}
(Spearman\,$\rho=0.44$, $n=9{,}905$) and with Thorsson immune subtypes
C1--C6 (Cramér's\,$V=0.44$). Because the tertiles are pan-cancer, the
proportion of hot tumours varied markedly across cancer types, from
$\leq1\%$ in kidney chromophobe, glioblastoma, pheochromocytoma and
paraganglioma, and brain lower-grade glioma to $>60\%$ in cervical,
head-and-neck squamous cell carcinoma and testicular germ cell tumours
(bladder cancer: 36\,\% hot). The immune and senescence landscape is
shown as a UMAP embedding in Fig.~\ref{fig:umap}.

\subsection{CDKN2A expression correlates positively with immune
infiltration (H1)}

Pan-cancer, CDKN2A expression and IFN-$\gamma$ score showed a
positive correlation (pooled Spearman $\rho=0.158$; median within-type
$\rho=0.112$; Fig.~\ref{fig:scatter}), significant at the primary
threshold ($\alpha=0.01$). This was a single Spearman test on the pooled
cohort; BH-FDR was applied only to the exploratory per-type and
correlation-table analyses. The within-type correlation was positive in
most but not all tumour types and heterogeneous in size: it was
strongest in kidney chromophobe ($\rho=0.44$), thyroid carcinoma
(0.40) and ovarian serous carcinoma ($>0.35$), and close to zero in
LGG (0.02), PAAD (0.04) and several other types.
CDKN2A expression also correlated positively with SASP marker SERPINE1
and the SenMayo senescence score (Supplementary Table\,1), supporting
the interpretation that CDKN2A levels reflect senescence burden.
\textbf{H1 is supported.}

\subsection{SASP scores are elevated in immune-hot tumours (H2)}

The custom SASP ssGSEA score differed across the three immune groups
(three-group Kruskal-Wallis, $p<0.001$; median score highest in hot and
lowest in cold tumours; Fig.~\ref{fig:alterations}b). This pooled
test does not account for tumour type. The Kruskal-Wallis test was run
for the SASP score only: the other senescence gene sets
(FRIDMAN\_SENESCENCE\_UP, REACTOME\_CELLULAR\_SENESCENCE,
SenMayo\cite{Saul2022}) were examined by Spearman correlation with the
IFN-$\gamma$ and CD8 scores, not by hot/cold group test. No analysis of
how the SASP difference varies across tumour types, and no per-gene
fold-change analysis (for example of SERPINE1), was performed.
\textbf{H2 is supported} for the SASP score.

% ── Figure 1 ───────────────────────────────────────────────────────────────
\begin{figure*}[t]
  \centering
  \includegraphics[width=\textwidth]{fig1_senescence_heatmap.png}
  \caption{\textbf{Pan-cancer senescence and immune phenotype
  co-enrichment.}
  Heatmap of the per-tumour-type median of seven scores (\textit{CDKN2A}
  expression, SASP, senescence effectors, FRIDMAN\_SENESCENCE\_UP,
  HALLMARK interferon-$\gamma$ response, SenMayo, HALLMARK p53 pathway),
  each z-scored across the 32 tumour types; cell labels give the raw
  median. Rows are tumour types and columns are scores. The figure is
  not stratified by hot/cold phenotype and does not itself test
  co-enrichment of senescence and immune phenotype.
  $n=9{,}910$ samples.}
  \label{fig:heatmap}
\end{figure*}

\subsection{CDKN2A genomic deletion is enriched in cold tumours after stratifying by tumour type (H3)}

Using cBioPortal alteration calls from the TCGA PanCancer Atlas 2018,
shallow and deep deletions of \textit{CDKN2A} were more frequent in hot
(38.6\,\%) than cold (31.7\,\%) tumours in the pooled comparison ($\chi^2$,
$p<0.001$; Fig.~\ref{fig:alterations}), which is confounded by tumour
type. Stratifying by tumour type reversed the direction (Mantel-Haenszel
OR\,=\,1.30, 95\,\%\,CI\,1.12--1.50, $p=4\times10^{-4}$; cold\,$>$\,hot in
21 of 29 tumour types with both phenotypes and both deletion states
represented). \textit{CDKN2B} deletion was, like CDKN2A, more frequent
in hot than cold tumours in the pooled comparison (38.2\,\% vs
31.5\,\%); no tumour-type-stratified test was run for CDKN2B, so no
stratified claim is made for it. \textit{TP53} mutation was more
prevalent in hot tumours in the pooled comparison (not stratified). \textbf{H3 is supported only in the tumour-type-stratified analysis.}

% ── Figure 2 ───────────────────────────────────────────────────────────────
\begin{figure*}[t]
  \centering
  \begin{subfigure}[t]{0.47\textwidth}
    \includegraphics[width=\linewidth]{alteration_hot_cold_grouped_bar.png}
    \caption{}
  \end{subfigure}
  \hfill
  \begin{subfigure}[t]{0.51\textwidth}
    \includegraphics[width=\linewidth]{sasp_score_hot_cold.png}
    \caption{}
  \end{subfigure}
  \caption{\textbf{Genomic alterations and SASP enrichment stratified
  by immune phenotype.}
  \textbf{(a)} Percentage of samples altered in hot, intermediate, and
  cold tumours pan-cancer for eight genes: CDKN2A, TP53, RB1, ATM and
  PTEN (deletion and/or mutation) and CDKN1A, CDKN2B and CDKN2C
  (deletion only). Bars show point estimates without confidence
  intervals; asterisks mark the pooled three-group $\chi^2$ test
  ($^*p<0.05$, $^{**}p<0.01$). The pooled CDKN2A deletion frequency is
  higher in hot (38.6\,\%) than cold (31.7\,\%) tumours; the
  direction reverses only after stratifying by tumour type (see text).
  \textbf{(b)} Pan-cancer SASP ssGSEA score distributions by immune
  phenotype, shown as overlaid histograms. SASP scores differ across
  groups (Kruskal-Wallis $p<0.001$).}
  \label{fig:alterations}
\end{figure*}

\subsection{Senescence and immune phenotype versus overall
survival (H4)}

The analysis code classifies samples into six groups (IFN-$\gamma$
tertile\,$\times$\,p16-high/low) using within-type median splits for CDKN2A
expression, and retains only the six-group cells with $\geq30$ samples
and $\geq5$ events (Fig.~\ref{fig:survival}). The global log-rank test
without stratification gave $p=3.3\times10^{-3}$, but stratifying by tumour
type gives $p=0.08$. The pre-specified ordering was not observed (median OS:
cold/p16-low 9.5\,y, hot/p16-high 7.8\,y). \textbf{H4 is not supported.}

\subsection{Immune and senescence markers as independent survival
predictors (H5)}

% ── Table 1 ────────────────────────────────────────────────────────────────
\begin{table*}[t]
\centering
\caption{\textbf{Cox proportional hazards regression results.}
OS model is stratified by tumour type and adjusted for age and
log(TMB\,+\,1) as well as the markers shown; it is \textbf{not}
purity-adjusted (purity data were unavailable). The competing-risks
model (Fine-Gray; not stratified by tumour type) uses
cancer-specific death as the primary event. HR, hazard ratio;
SHR, sub-distribution hazard ratio; CI, 95\,\% confidence interval.
OS Cox model: $n\approx9{,}827$ samples, $\approx2{,}748$ OS events (of
9,850 samples with OS data and 2,751 events).}
\label{tab:cox}
\renewcommand{\arraystretch}{1.3}
\fontsize{8}{10}\selectfont
\begin{tabularx}{\textwidth}{%
  l
  S[table-format=1.3] @{\ } >{\footnotesize}l @{\ } S[table-format=1.2e1]
  S[table-format=1.3] @{\ } >{\footnotesize}l @{\ } S[table-format=1.2e1]}
\toprule
& \multicolumn{3}{c}{\textbf{Overall Survival (Cox PH)}}
& \multicolumn{3}{c}{\textbf{Cancer-Specific Death (Fine-Gray)}} \\
\cmidrule(lr){2-4}\cmidrule(lr){5-7}
\textbf{Variable}
  & {\textbf{HR}} & \textbf{95\,\% CI} & {$p$-value}
  & {\textbf{SHR}} & \textbf{95\,\% CI} & {$p$-value} \\
\midrule
Age (years)
  & 1.026 & (1.022--1.030) & {$2.1\times10^{-46}$}
  & 1.017 & (1.013--1.021) & {$<10^{-10}$} \\
IFN-$\gamma$ score
  & 0.907 & (0.856--0.961) & {$8.5\times10^{-4}$}
  & 0.749 & (0.704--0.797) & {$<10^{-10}$} \\
\textit{CDKN2A} expression
  & 0.989 & (0.973--1.005) & {0.160\ \ \ \ }
  & 1.050 & (1.028--1.072) & {$6.4\times10^{-6}$} \\
\textit{CDKN1A} expression (p21)
  & 0.935 & (0.897--0.975) & {$1.6\times10^{-3}$}
  & 0.887 & (0.850--0.925) & {$3.2\times10^{-8}$} \\
\textit{SERPINE1} expression (PAI-1)
  & 1.161 & (1.133--1.190) & {$1.5\times10^{-33}$}
  & 1.245 & (1.212--1.278) & {$<10^{-10}$} \\
\textit{GDF15} expression
  & 1.002 & (0.980--1.024) & {0.892\ \ \ \ }
  & 0.948 & (0.932--0.965) & {$1.2\times10^{-9}$} \\
log(TMB)
  & 0.922 & (0.874--0.973) & {$3.1\times10^{-3}$}
  & 1.044 & (0.990--1.100) & {0.110\ \ \ \ } \\
\bottomrule
\end{tabularx}
\end{table*}

IFN-$\gamma$ score was a highly significant independent predictor of OS
(HR\,=\,0.907, 95\,\%\,CI\,0.856--0.961, $p=8.5\times10^{-4}$;
Table\,\ref{tab:cox}; Fig.~\ref{fig:forest}a). SERPINE1/PAI-1 was the
strongest individual risk factor (HR\,=\,1.161, $p=1.5\times10^{-33}$).
p21/CDKN1A was independently protective (HR\,=\,0.935, $p=1.6\times10^{-3}$),
as was higher TMB (HR\,=\,0.922, $p=0.003$). \textit{CDKN2A} expression
itself did not reach significance in the OS model (HR\,=\,0.989,
$p=0.160$), failing to confirm H5 at the Bonferroni threshold. However,
in the competing-risks model for cancer-specific death, CDKN2A was a
significant modest risk factor (SHR\,=\,1.050, $p=6.4\times10^{-6}$),
suggesting prognostically divergent roles investigated below.
\textbf{H5 is not confirmed} at the pre-specified Bonferroni
$\alpha=0.01$.

\subsection{A significant CDKN2A\,$\times$\,IFN-$\gamma$ interaction
modifies survival}

The opposing directions of the CDKN2A coefficient in OS versus
cancer-specific death, combined with its positive correlation with
IFN-$\gamma$, suggested effect modification. An interaction Cox model
including CDKN2A\_expr\,$\times$\,IFN-$\gamma$\_score yielded a highly
significant interaction term (HR\,=\,0.964 per unit increase in the
product, 95\,\%\,CI\,0.949--0.980, $p=1.7\times10^{-5}$;
Fig.~\ref{fig:forest}b). This exploratory (not pre-specified) result is
consistent with CDKN2A expression being protective mainly in the
presence of IFN-$\gamma$ signalling, but it is observational and
unreplicated.

% ── Figure 3 ───────────────────────────────────────────────────────────────
\begin{figure*}[t]
  \centering
  \begin{subfigure}[t]{0.48\textwidth}
    \includegraphics[width=\linewidth]{cdkn2a_vs_ifng_scatter.png}
    \caption{}
  \end{subfigure}
  \hfill
  \begin{subfigure}[t]{0.48\textwidth}
    \includegraphics[width=\linewidth]{cdkn2a_ifng_corr_by_type.png}
    \caption{}
  \end{subfigure}
  \caption{\textbf{CDKN2A expression positively correlates with
  IFN-$\gamma$ score pan-cancer (H1).}
  \textbf{(a)} Pan-cancer scatter plot of \textit{CDKN2A}
  log$_2$(norm\_count\,+\,1) versus IFN-$\gamma$ score ($n=9{,}910$);
  points are a single colour and no fit line is drawn. Pooled Spearman
  $\rho=0.158$.
  \textbf{(b)} Per-tumour-type Spearman $\rho$(\textit{CDKN2A},
  IFN-$\gamma$) shown as a ranked bar chart (tumour types with
  $\geq30$ samples; no significance markers are drawn). Most tumour
  types show a positive correlation (median $\rho=0.112$).}
  \label{fig:scatter}
\end{figure*}

\subsection{Tumour-type heterogeneity in immune-survival coupling}

Per-tumour-type log-rank tests (three groups: cold, intermediate, hot;
tumour types with $n\geq50$; BH-FDR) were significant at FDR\,$<$\,0.05
in three cancer types: LGG ($p_\text{adj}=2.3\times10^{-14}$), KIRP
($p_\text{adj}=3.1\times10^{-5}$) and PAAD ($p_\text{adj}=0.017$).
LUAD ($p_\text{adj}=0.0501$) and UVM ($p_\text{adj}=0.060$) narrowly
missed the threshold and are not significant. These are non-directional
three-group tests, so they show that survival differs across immune
groups but not that hot tumours fare better. LGG gave the smallest
$p$-value (raw $p=7.8\times10^{-16}$) but contains only 4 hot tumours
(24 intermediate, 486 cold), so this cannot be read as a hot-tumour
effect and may reflect the confounding of IDH mutation status with
immune phenotype. Bladder and cervical cancers showed no significant
difference ($p_\text{adj}=0.24$ for both).

% ── Figure 4: Survival ─────────────────────────────────────────────────────
\begin{figure*}[t]
  \centering
  \includegraphics[width=\textwidth]{fig4_survival_composite.png}
  \caption{\textbf{Kaplan-Meier survival curves (unstratified; see text).}
  Montage of the Kaplan-Meier figures produced by the survival analysis
  (every \texttt{km\_*.png} file, each panel titled by its file name):
  overall survival by cold/intermediate/hot immune phenotype; overall
  survival by the six immune-tertile\,$\times$\,p16-high/low groups
  (file \texttt{km\_4group.png}, a legacy name; unstratified log-rank
  $p=3.3\times10^{-3}$, stratified by tumour type $p=0.077$); TP53
  mutated versus wild-type; disease-specific survival and
  progression-free interval by immune phenotype; and overall survival by
  within-type median age. The curves are not adjusted for tumour type,
  age or purity. The forest plots are shown in Fig.~\ref{fig:forest}.}
  \label{fig:survival}
\end{figure*}

% ── Figure 5: Forest plots ─────────────────────────────────────────────────
\begin{figure*}[t]
  \centering
  \begin{subfigure}[t]{0.48\textwidth}
    \includegraphics[width=\linewidth]{cox_forest_plot.png}
    \caption{}
  \end{subfigure}
  \hfill
  \begin{subfigure}[t]{0.48\textwidth}
    \includegraphics[width=\linewidth]{competing_risks_forest.png}
    \caption{}
  \end{subfigure}
  \caption{\textbf{Cox proportional hazards and competing-risks forest
  plots.}
  \textbf{(a)} Forest plot from the main Cox PH model for overall
  survival (strata\,=\,tumour type; $n\approx9{,}827$; $\approx2{,}748$ OS events).
  Points represent HR; bars represent 95\,\% CI; dashed line at HR\,=\,1.
  The model is not purity-adjusted. IFN-$\gamma$ score (HR\,=\,0.907) and TMB
  (HR\,=\,0.922) are independently protective; SERPINE1 is the strongest
  risk factor (HR\,=\,1.161).
  \textbf{(b)} Competing-risks (Fine-Gray) forest plot for cancer-specific
  death. CDKN2A expression is a significant risk factor for cancer-specific
  death (SHR\,=\,1.050, $p=6.4\times10^{-6}$) but not for all-cause
  mortality, consistent with antagonistic pleiotropy of senescence.}
  \label{fig:forest}
\end{figure*}

% ── Figure 6: UMAP ─────────────────────────────────────────────────────────
\begin{figure*}[t]
  \centering
  \includegraphics[width=\textwidth]{fig2_umap.png}
  \caption{\textbf{UMAP embedding of the TCGA pan-cancer immune and
  senescence landscape.}
  UMAP of 9,905 TCGA samples projected on five features (IFN-$\gamma$
  score, \textit{CDKN2A} expression, SenMayo NES, SASP score, CD8$^+$
  fraction; umap-learn v0.5; \textit{n}\_neighbors\,=\,30,
  min\_dist\,=\,0.3, random\_state\,=\,42). Panels are coloured by
  \textbf{(left)} immune phenotype: hot (crimson, \#B2182B),
  intermediate (grey, \#878787), cold (deep blue, \#2166AC);
  \textbf{(centre)} \textit{CDKN2A} genomic alteration status
  (deleted, orange; intact, grey); and
  \textbf{(right)} SenMayo senescence score (RdBu\textsubscript{r}
  continuous scale). Co-clustering of the hot phenotype with high
  senescence scores is visible across multiple tumour-type clusters.}
  \label{fig:umap}
\end{figure*}

% ═══════════════════════════════════════════════════════════════════════════
\section{Discussion}

Our analysis provides pan-cancer observational associations between
p16/CDKN2A expression, senescence scores and immune phenotype; it cannot
establish that senescence shapes the microenvironment or patient survival. Of five pre-specified confirmatory
hypotheses, H1 and H2 were supported, H3 only after stratification by
tumour type, and H4 was not supported; H5 (CDKN2A as
independent OS predictor) failed significance, with the interaction
analysis giving a more nuanced (exploratory) picture.

The positive correlation between CDKN2A expression and IFN-$\gamma$
score (H1) is consistent with mechanistic studies showing that SASP
cytokines — particularly CXCL9, CXCL10, and IL-6 — directly recruit
CD8$^+$ T cells and NK cells to tumour sites\cite{Krtolica2001}. Our
finding that CDKN2A deletion is enriched in cold tumours only after
stratifying by tumour type (H3; the pooled comparison points the other
way) may add a genomic dimension: cancers that have lost CDKN2A could be
cold because they have also lost the capacity to mount a
senescence-driven immune recruitment programme, though this is
speculative. Therapeutically, SASP restoration through
CDK4/6 inhibition — which can induce senescence in neighbouring
cells\cite{Goel2017} — or epigenetic CDKN2A re-expression might shift
cold tumours toward an immune-permissive state.

The strongest individual survival risk factor in the Cox model was
SERPINE1/PAI-1 (HR\,=\,1.161 for OS; SHR\,=\,1.245 for cancer death),
a pleiotropic SASP component that promotes extracellular matrix
remodelling, angiogenesis, and fibrin-mediated immune
evasion\cite{Alessi1997}. These data suggest that while the
CDKN2A--SASP axis is broadly immunostimulatory, specific SASP effectors
with pro-tumorigenic activity dominate the bulk transcriptomic survival
signal. Deconvolving the ``helpful'' (CXCL9/10, IL-15) from the
``harmful'' (SERPINE1, TGF-$\beta$) SASP in different tumour contexts is
an important open challenge.

The failure of CDKN2A expression to independently predict OS (H5) is
interpretable in light of the significant CDKN2A\,$\times$\,IFN-$\gamma$
interaction (HR\,=\,0.964, $p=1.7\times10^{-5}$): CDKN2A is only
protective in IFN-$\gamma$-high tumours. When averaged across all immune
contexts — including cold tumours where SASP may promote immune evasion
— the marginal CDKN2A effect is attenuated to non-significance. The
divergence of the CDKN2A coefficient between OS (HR\,$\approx$\,1.0)
and cancer-specific death (SHR\,=\,1.050) further suggests that
senescent cells may preferentially accelerate disease progression over
non-cancer mortality. This mirrors the canonical dual role of senescence:
cell-autonomous tumour suppression versus non-cell-autonomous SASP
promotion of the tumour microenvironment\cite{Krtolica2001}.

Several limitations must be acknowledged. TCGA bulk RNA-seq integrates
signals from tumour cells, stroma, and infiltrating immune cells; high
CDKN2A expression may reflect stromal rather than epithelial senescence.
Tumour-purity estimates could not be obtained (the download failed),
so no analysis was purity-adjusted and stromal confounding is not
addressed. ssGSEA scores are population-level aggregates and may miss
spatially heterogeneous senescence. Causal inference is limited;
reverse causation — hot tumours inducing paracrine senescence through
IFN-$\gamma$ or cytotoxicity\cite{Acosta2013} — cannot be excluded.
Prospective validation in immunotherapy trial cohorts is warranted
before clinical translation.

Despite these limitations, this study establishes a foundation for
mechanistic investigation of the CDKN2A--SASP--immunity axis and for
the rational combination of senescence-inducing therapies with immune
checkpoint blockade, particularly in cold tumours, where CDKN2A
deletion was more frequent after stratification by tumour type.

% ═══════════════════════════════════════════════════════════════════════════
\section{Methods}

\subsection{Data acquisition}
RNA-seq data (gene-level RSEM, log$_2$(norm\_count\,+\,1);
\texttt{EB++AdjustPANCAN\_IlluminaHiSeq\_RNASeqV2}) were downloaded from
the UCSC Xena public hub (TCGA Pan-Cancer; not the GDC pipeline). The
analysis retains primary-tumour samples only. The pipeline attempts to
download consensus tumour-purity estimates (Aran
\textit{et al.}\cite{Aran2015}) from Xena, but the request was refused
(HTTP 403), so no purity data were obtained and no analysis uses purity.
Thorsson immune subtypes and CIBERSORT fractions were from the TCGA
Immune Landscape\cite{Thorsson2018}. CDKN2A alteration calls were
retrieved via the cBioPortal REST API
(\texttt{tcga\_pan\_can\_atlas\_2018}). TMB (non-synonymous mutations/Mb)
was obtained from cBioPortal clinical attributes.

\subsection{Immune phenotype classification}
IFN-$\gamma$ scores were computed as the mean of per-gene z-scores
(each gene z-scored across all samples) over \textit{IFNG},
\textit{STAT1}, \textit{CCL5}, \textit{CXCL9}, \textit{CXCL10}, and
\textit{IDO1}. Samples were classified as hot/intermediate/cold by
pan-cancer tertile (not within tumour type), so the hot fraction differs
strongly between tumour types and pooled comparisons can be confounded
by tumour type.

\subsection{Senescence scoring}
ssGSEA scores were computed using gseapy\,(v1.1)\cite{Fang2023} with:
\texttt{FRIDMAN\_SENESCENCE\_UP} and \texttt{REACTOME\_CELLULAR\_SENESCENCE}
from MSigDB v2023.1; a custom 16-gene SASP list; and the SenMayo
signature\cite{Saul2022} (\texttt{SAUL\_SEN\_MAYO}; 124 genes in the cached
MSigDB file); further sets (e.g.\ GOBP and Hallmark sets) were scored but
are not reported here. The Fig.~\ref{fig:heatmap} heatmap z-scores
per-tumour-type medians across tumour types. The ssGSEA computation was cached to
disk.

\subsection{UMAP}
UMAP was computed on five immune and senescence features (IFN-$\gamma$
score, \textit{CDKN2A} expression, SenMayo NES, SASP score, CD8$^+$
fraction) across 9,905 samples with complete feature data, using
umap-learn\,(v0.5) with \texttt{n\_neighbors=30}, \texttt{min\_dist=0.3},
\texttt{random\_state=42}.

\subsection{Statistical testing}
Spearman correlations were computed per tumour type (types with
$n\geq30$) and pan-cancer; the H1 pan-cancer test is a single Spearman
test judged against the primary threshold, and BH-FDR correction was
applied to the exploratory per-type and correlation-table analyses. A
three-group Kruskal-Wallis test compared the SASP score by immune
phenotype. Chi-squared tests compared CDKN2A alteration frequencies
across the immune groups in the full cohort and in the subset of tumour
types with $n\geq100$. Because the pooled comparison is confounded by
tumour type, a Mantel-Haenszel test (cold versus hot, strata = tumour
type, restricted to strata containing both phenotypes and both deletion
states) was used to judge H3.

\subsection{Survival analysis}
Kaplan-Meier curves and log-rank tests were computed using the
\texttt{survival}\,(v3.6) package in R. The six-group log-rank test
(immune tertile\,$\times$\,within-type p16 median split; cells with
$\geq30$ samples and $\geq5$ events) was reported both unstratified and
stratified by tumour type. Cox PH models used \texttt{coxph()}
stratified on tumour type. Covariates: age, IFN-$\gamma$ score, CDKN2A
expression, CDKN1A, SERPINE1, GDF15 expression and log(TMB\,+\,1);
tumour purity was not available and is not in any model. The
CDKN2A\,$\times$\,IFN-$\gamma$ interaction model was an exploratory
analysis. The Fine-Gray competing-risks model used \texttt{crr()} from
the \texttt{cmprsk} package (no tumour-type strata). Per-tumour-type
log-rank tests (three immune groups, non-directional) were run in tumour
types with $n\geq50$. BH-FDR correction was applied separately to OS and DSS per-type
$p$-values.

\subsection{Multiple testing}
H1--H5 used Bonferroni correction ($\alpha=0.01$ across five primary
tests). All other tests are exploratory, reported with BH-FDR
$q$-values at $\alpha=0.05$.

\subsection{Code and data availability}
All analysis code is available at [GitHub URL]. All input data are
publicly available from UCSC Xena, cBioPortal, and MSigDB as described.
No new data were generated.

% ═══════════════════════════════════════════════════════════════════════════
\section*{Acknowledgements}
We acknowledge the TCGA Research Network
(\url{https://www.cancer.gov/tcga}) for generating the data used in
this analysis and the UCSC Xena team for data hosting infrastructure.

\section*{Author contributions}
D.C.K. conceived the study, performed all analyses, and wrote the
manuscript.

\section*{Competing interests}
The authors declare no competing interests.

% ── References ─────────────────────────────────────────────────────────────
\begin{thebibliography}{99}
\fontsize{7.5}{10}\selectfont

\bibitem{Chen2018}
Chen, D.\,S. \& Mellman, I. Elements of cancer immunity and the
cancer--immune set point. \textit{Nature} \textbf{541}, 321--330 (2018).

\bibitem{Sharma2017}
Sharma, P. \& Allison, J.\,P. The future of immune checkpoint therapy.
\textit{Science} \textbf{348}, 56--61 (2017).

\bibitem{Coppé2010}
Copp{\'e}, J.-P. \textit{et al.} The senescence-associated secretory
phenotype: the dark side of tumour suppression.
\textit{Annu. Rev. Pathol.} \textbf{5}, 99--118 (2010).

\bibitem{Acosta2013}
Acosta, J.\,C. \textit{et al.} A complex secretory program orchestrated
by the inflammasome controls paracrine senescence.
\textit{Nature Cell Biol.} \textbf{15}, 978--990 (2013).

\bibitem{Kuilman2008}
Kuilman, T. \textit{et al.} Oncogene-induced senescence relayed by an
interleukin-dependent inflammatory network.
\textit{Cell} \textbf{133}, 1019--1031 (2008).

\bibitem{Rayess2012}
Rayess, H., Wang, M.\,B. \& Srivatsan, E.\,S. Cellular senescence and
tumour suppressor gene p16.
\textit{Int. J. Cancer} \textbf{130}, 1715--1725 (2012).

\bibitem{Thorsson2018}
Thorsson, V. \textit{et al.} The immune landscape of cancer.
\textit{Immunity} \textbf{48}, 812--830 (2018).

\bibitem{Saul2022}
Saul, D. \textit{et al.} A new gene set identifies senescent cells and
predicts senescence-associated pathways across tissues.
\textit{Nature Aging} \textbf{2}, 1187--1202 (2022).

\bibitem{Goel2017}
Goel, S. \textit{et al.} CDK4/6 inhibition triggers anti-tumour
immunity. \textit{Nature} \textbf{548}, 471--475 (2017).

\bibitem{Alessi1997}
Alessi, M.\,C. \textit{et al.} Plasminogen activator inhibitor 1,
transforming growth factor-$\beta_1$, and risk of cardiovascular events.
\textit{Circulation} \textbf{96}, 2252--2258 (1997).

\bibitem{Krtolica2001}
Krtolica, A. \textit{et al.} Senescent fibroblasts promote epithelial
cell growth and tumorigenesis.
\textit{Proc. Natl Acad. Sci. USA} \textbf{98}, 12072--12077 (2001).

\bibitem{Aran2015}
Aran, D., Sirota, M. \& Butte, A.\,J. Systematic pan-cancer analysis of
tumour purity. \textit{Nature Commun.} \textbf{6}, 8971 (2015).

\bibitem{Fang2023}
Fang, Z. \textit{et al.} GSEApy: a comprehensive package for performing
gene set enrichment analysis in Python.
\textit{Bioinformatics} \textbf{39}, btac757 (2023).

\end{thebibliography}
\end{document}
